

%\begin{lemma}
%	For any profinite set $S$ there is a natural diagram $D_S\colon I_S\to \Top$ of finite discrete spaces with a natural isomorphism
%		$$S\to\limit D_S =\limit_{i\in I_S} D_S(i),$$
%	presenting $S$ as a projective limit of finite discrete spaces. 
%\end{lemma}
%\begin{proof}
%	For each profinite $S$ consider the category $I_S = S/\N \subseteq S/\Top$, the under-category of finite, discrete spaces, where we identify $n\in \N$ with the discrete space $\{0, \dots, n\}$. Consider the diagram $D_S\colon I_S\to \Top$ given by mapping $S\to S'$ to the finite discrete space $S'$. This construction is clearly natural in $S$.
%	
%	Now all the $S\to S'$ in $D_S$, already constitute a cone over $D_S$ with tip $S$, hence there is a natural morphism $S\to \limit D_S$. Again this construction (by functoriality of taking limits) is natural in $S$.
%	
%	It remains to show that this morphism is an isomorphism. Clearly since $S$ is profinite there is a presentation $S\cong \limit_{i\in I} S_i$ of $S$ by finite discrete spaces $S_i$. We thus have induced morphisms $p_i\colon S\to \{0, \dots, n_i\}$, where $S_i \cong \{0, \dots, n_i\}$.
%	
%	\begin{todo}
%		finish this
%	\end{todo}
%\end{proof}

\begin{definition}[Final topologies]
	\label{definition: final topologies}
	
	Let $X$ be a set and $(f_i\colon X_i \to X)_{i\in I}$ a family of functions from topological spaces $X_i$ into $X$. The finest topology on $X$ such that all maps $f_i$ are continuous is called the \emph{final topology}. Equivalently, the final topology is the collection of sets $U\subseteq X$ such that $f_i^{-1}(U)$ is open in $X_i$ for any $i\in I$. That is a subset of $X$ is open if and only if its preimage under any $f_i$ is open.
\end{definition}

%\begin{remark}
%	The inclusion $i$ of $\kappa$-compactly generated spaces into $\Top$ has a right adjoint $X\mapsto X^{\boundedkificationsymbol \kappa}$ where $X^{\boundedkificationsymbol \kappa}$ has the same underlying set as $X$ and carries the final topology induced by all continuous maps $S\to X$ from a $\kappa$-small compact Hausdorff space $S$. The unit and counit of the adjunction $i\ladj (-)^{\boundedkificationsymbol{\kappa}}$ are both induced by the identities of the underlying sets. 
%	
%	\begin{todo}
%		%
%		proof or reference
%	\end{todo}
%\end{remark}

\begin{definition}[Ultrafilters, principal ultrafilters \& the Stone-topology]
	\label{definition: ultrafilters, principal ultrafilters and the stone-topology}
	
	Let $X$ be a set. A $\emptyset \ne F \subseteq \powerset X$ is called an \emph{ultrafilter} on $X$ if
	\begin{enumerate}
		\item $\emptyset \not\in F$.
		\item If $U \in F$ and $U\subseteq V\subseteq X$ then $V \in F$.
		\item If $U, V\in F$ then $U\cap V\in F$.
		\item If $U\subseteq X$ either $U \in F$ or $X\setminus U\in F$.
	\end{enumerate}
	That is, an ultrafilter is a maximal, non-trivial filter on $X$. Write
		$$\filterkernel F \ldef \bigcap_{U\in F} U.$$
	If $F$ is an ultrafilter on $X$ and in addition
		$$\filterkernel F = \{x\}$$
	for some $x \in X$, then $F$ is called the \emph{principal ultrafilter} at $x$ and $F$ is uniquely determined by $x$. The set of ultrafilters on $X$ admits a topology generated by the sets of ultrafilters $\{F \setseparator U \in F\}$ where $U \subseteq X$. This topology is called the \emph{Stone} topology.
\end{definition}

%\begin{proposition}[Stone-\v Cech compactification of discrete spaces]
%	\label{proposition: Stone-Cech of discrete spaces}
%	If $X$ is a discrete topological space, then the Stone-\v Cech compactification of $X$ is an extremally disconnected profinite set and is given by the set of all ultrafilters on $X$ together with the Stone topology. The map $X\to\stonecech X$ is given by mapping each $x\in X$ to the principal ultrafilter at $x$.
%\end{proposition}
%
%\begin{proof}
%	%Give reference
%%	We verify the universal property. To that end let $K$ be a compact Hausdorff space and $f\colon X \to K$ a continuous map. We construct a map $\stonecech f\colon \stonecech X\to K$.
%%	
%%	We need the following result:
%%	Consider any ultrafilter $F'$ on $K$. Suppose we have distinct points $k, k' \in \filterkernel F'$. As $K$ is Hausdorff we have neighborhoods $U$ of $k$ and $V$ of $k'$ such that $U\cap V=\emptyset$. As $F'$ is an ultrafilter either $U \in F'$ or $K\setminus U \in F'$. Suppose $K\setminus U \in F'$, then $x \in \filterkernel F' \subseteq \bigcap_{V \in F'} V \subseteq K\setminus U$, but $x\in U$. As this is not possible, we have $U \in F'$. Similarly $V \in F'$ and hence $\emptyset = U\cap V \in F'$. A contradiction.
%%	
%%	Now consider some ultrafilter $F \in \stonecech X$ and set
%%	$$\{(\stonecech f)(F)\} \ldef \bigcap_{U \in F} f(U).$$
%%	Indeed this is well-defined, since
%%	$$\bigcap_{U \in F} f(U) = \bigcap_{V \in f(F)} V = \filterkernel (f(F)^\uparrow)$$
%%	where $f(F) = \{f(U) \setseparator U \in F\}$ and $(-)^\uparrow$ denotes upward closure, making $f(F)$ into an ultrafilter on $K$. Hence $\bigcap_{U \in F} f(U) = \filterkernel f(F)^\uparrow$ is a singleton.
%%	
%%	%: define upward closure
%%	
%%	Now if $i\colon X \to \stonecech X$ denotes the inclusion that sends each $x\in X$ to the principal ultrafilter $F_x = \{U \setseparator x\in U\}$ at $x$, then $f(x) \in \filterkernel f(F_x)^\uparrow$, as any set of the form $f(U)$ with $x\in U$ contains $f(x)$. Hence $(\stonecech f)(F_x) = f(x)$ for all $x \in X$, so $\stonecech f \circ i = f$ and the triangle commutes.
%%	
%%	It remains to show that $\stonecech f$ is continuous.  %
%%	
%\end{proof}


%\begin{lemma}
%	If $X\colon \category I \to \Top$ is a diagram of compact Hausdorff spaces, then $\limit X$ is compact.
%\end{lemma}
%
%\begin{proof}
%	Embed $\limit X \subseteq \prod_{i\in \category I} X_i$. Since all $X_i$ are compact, so is their product. It remains to show that $\limit X$ is a closed subset of the product, hence compact.
%	
%	For each $f\colon i\to j$ in $\category I$ denote by $\Gamma_f \subseteq X_i \times X_j$ the graph of $X(f)$. Since $\Gamma_f$ is the inverse image of the diagonal under the natural map $X_i \times X_j \to X_j \times X_j$ and the diagonal is closed since $X_j$ is Hausdorff we find $\Gamma_f$ to be closed.
%	
%	% double use of index $i$ here
%	Now
%		$$\limit X = \{(x_i)_{i\in\category I}\setseparator\forall i\to j\colon X(i\to j)(x_i) = x_j \}=\bigcap_{f\colon i\to j}\{(x_i)_{i\in\category I}\setseparator (x_i, x_j) \in \Gamma_f\}$$
%	is the intersection of the Cartesian products of the $\Gamma_f$ with all other $X_k$ for $k\ne i, j$. These products are closed and hence so is $\limit X$ itself.
%\end{proof}
